\documentclass[10pt,a4paper]{article}
\usepackage[hmargin=1.5cm,top=1.1cm,bottom=1.1cm]{geometry}
\usepackage[T1]{fontenc}
\usepackage[utf8]{inputenc}
\usepackage{lmodern}
\usepackage{microtype}
\usepackage{booktabs}
\usepackage{graphicx}
\usepackage{xcolor}
\usepackage{tcolorbox}
\usepackage{enumitem}
\usepackage[hidelinks]{hyperref}
\usepackage{soul}
\sethlcolor{yellow!55}
\soulregister\texttt7
\soulregister\emph7
\soulregister\textbf7

\setlist{nosep,leftmargin=*}
\setlength{\parskip}{2.5pt}
\setlength{\parindent}{0pt}
\pagestyle{empty}
\newcommand{\figdir}{../output/figures}
\newcommand{\q}[1]{\textbf{\color{blue!50!black}#1}}

\begin{document}
\small

{\large\textbf{Does respecting the image pay? Convolutional networks on Fashion-MNIST}}\\
Ignacio Queipo de Llano \textperiodcentered{} Predictive Machine Learning in Finance, IESE MiF \textperiodcentered{} 27 November 2026\\[-6pt]
\rule{\linewidth}{0.4pt}

\begin{tcolorbox}[colback=white,colframe=black,boxrule=0.4pt,sharp corners,left=5pt,right=5pt,top=3pt,bottom=3pt]
\textbf{Bottom line.} Keeping the 2-D structure pays: \textbf{the case's CNN reaches 90.0\% in 5 epochs against
86.3\% for a dense network}, over a quarter fewer errors, and a deeper CNN with batch normalisation and dropout
reaches \textbf{92.6\%}. The price is compute (125 s and 810 s on CPU against 43 s). What is left is mostly shirts,
63\% right, confused with T-shirts, pullovers and coats that a person would also hesitate over: the ceiling now is
the 28 $\times$ 28 grey image, not the model.
\end{tcolorbox}

\begin{minipage}[t]{0.50\linewidth}\vspace{0pt}
\q{Q1. The data.} Same images as this morning: (60000, 28, 28) and (10000, 28, 28), pixels 0--255 (divided by 255).
\hl{Tops (T-shirt, shirt, pullover, coat) are the hard ones.}

\q{Q2. Dense baseline} 784 $\to$ 32 $\to$ 32 $\to$ 10, 10 epochs: \hl{test 86.3\%}, 26,506 parameters (the same as
the scikit-learn (32, 32), which reached 87.2\% in 60 epochs).

\q{Q3. Running fit() again.}

\vspace{2pt}
\centerline{\begin{tabular}{lrr}
\toprule
Calls to fit() & Train & Test\\
\midrule
1 (epochs 1--10) & 89.3\% & 86.3\%\\
2 (epochs 11--20) & 90.5\% & 86.7\%\\
3 (epochs 21--30) & 91.2\% & 86.8\%\\
\bottomrule
\end{tabular}}

\vspace{2pt}
\hl{Keras keeps the weights and continues training; scikit-learn's fit() starts again} (unless \texttt{warm\_start}).
Train keeps rising and test stalls: overfitting builds up with each call.

\q{Q4. Shape (28, 28, 1).} A convolution needs the grid to know which pixels are neighbours; the last 1 is the
channel (grey; a colour photo would have 3).

\q{Q6. 3$\times$3 filters} learn edges (vertical, horizontal, diagonal), corners and textures (knit, stripes). Many
filters look for many patterns at once; the second layer combines them into parts (a collar, a sole).

\q{Q7. Max pooling} keeps the strongest response in each 2$\times$2 block, so a pattern shifted by a pixel gives the
same output: \hl{a sneaker is a sneaker wherever it sits in the frame}. It also quarters the grid.
\end{minipage}\hfill
\begin{minipage}[t]{0.47\linewidth}\vspace{0pt}
\q{Q5, Q8. The case's CNN} (\texttt{model.summary()}):

\vspace{2pt}
\centerline{\begin{tabular}{llr}
\toprule
Layer & Output & Parameters\\
\midrule
Conv2D 32, 3$\times$3 & 26$\times$26$\times$32 & 320\\
MaxPooling 2$\times$2 & 13$\times$13$\times$32 & 0\\
Conv2D 64, 3$\times$3 & 11$\times$11$\times$64 & 18,496\\
MaxPooling 2$\times$2 & 5$\times$5$\times$64 & 0\\
Flatten, Dense 64 & 64 & 102,464\\
Dense 10, softmax & 10 & 650\\
\midrule
Total & & \hl{\textbf{121,930}}\\
\bottomrule
\end{tabular}}

\vspace{2pt}
4.6 times the dense network. \hl{Weight sharing makes 32 filters cost 320 parameters; 84\% sit in the first dense
layer.}

\q{Q9. Five epochs.} Validation 87.3\% $\to$ 88.7 $\to$ 88.9 $\to$ 89.9 $\to$ 90.4\%; \hl{test 90.0\% vs 86.3\%
dense: +3.7 points, 27\% fewer errors}, in 125 s.

\includegraphics[width=\linewidth]{\figdir/q10_misclassified.png}

\q{Q10. Mistakes} (1,002). \hl{Shirt 62.9\%}, coat 81.8\%, T-shirt 86.6\%; sneaker 98.4\%, bag 97.8\%. Most errors
are tops at low resolution that a person would also hesitate over.
\end{minipage}

\vspace{3pt}
\q{Q11. Challenge: beat 90\%.}

\vspace{2pt}
\centerline{\begin{tabular}{lrrrrr}
\toprule
Model & Parameters & Epochs & Train & Test & CPU\\
\midrule
Dense 784-32-32-10 & 26,506 & 10 & 89.3\% & 86.3\% & 43 s\\
CNN of the case & 121,930 & 5 & 92.6\% & 90.0\% & 125 s\\
2 $\times$ [conv 3$\times$3, BN, conv 3$\times$3, BN, pool, dropout 0.25], dense 128, dropout 0.4 & 468,586 & 12 &
95.3\% & \hl{\textbf{92.6\%}} & 810 s\\
\bottomrule
\end{tabular}}

\vspace{2pt}
Two convolutions per stage see larger patterns before pooling; batch normalisation keeps each layer's inputs on a
stable scale (the morning's normalisation lesson, inside the network); dropout fights the overfitting of Q3.

\q{Q12. The journey.} \hl{Raw pixels failed because 0--255 inputs swamp small starting weights (31.6\%). Normalising
let the network learn (87\%). The CNN wins because it uses the image's structure: local filters, shared weights and
pooling (90.0\%, 92.6\% deeper).} The same idea, building the structure of the data into the model, gives sequence
models for time series and transformers for text.
\end{document}
